% Source PDF: source/普通高考/2015/2015北京文.pdf, page 1, question 5.
% The source PDF contains vector flowchart line art (no embedded bitmap); this
% redraw follows a no-grid render localized with page-1 grid evidence.
% Node-edge list: 开始→(k=0,a=3,q=1/2)→a=aq→k=k+1→a<1/4;
% 否→(left loop)→a=aq; 是→输出 k→结束. No dashed edges or other links.
% The left loop's horizontal arrow merges above a=aq; the main vertical arrow
% uniquely enters a=aq from the merge. Every node text is centered and zero-indented.

\begin{tikzpicture}[
  >=Stealth,
  line width=0.45pt,
  line cap=round,
  line join=round,
  font=\normalsize,
  flowtext/.style={anchor=center,align=center,inner sep=0pt,
    execute at begin node={
      \setlength{\parindent}{0pt}%
      \setlength{\leftskip}{0pt}%
      \setlength{\rightskip}{0pt}%
      \everypar{}%
    }},
  every node/.style={flowtext},
  flowbox/.style={flowtext,draw,minimum height=0.68cm,inner xsep=3pt,inner ysep=2pt},
  terminal/.style={flowbox,rounded corners=0.18cm,minimum width=1.05cm},
  io/.style={flowbox,shape=trapezium,trapezium left angle=72,trapezium right angle=72,
    minimum height=0.72cm},
  process/.style={flowbox,minimum width=1.55cm},
  decision/.style={flowbox,shape=diamond,aspect=2.65,minimum width=4.05cm,minimum height=1.02cm,inner sep=0pt},
]
  \node[terminal] (start) at (0,9.65) {开始};
  \node[process,minimum width=3.40cm,minimum height=0.82cm] (init) at (0,8.35)
    {$k=0,a=3,q=\dfrac12$};
  \node[process,minimum width=1.25cm] (multiply) at (0,6.95) {$a=aq$};
  \node[process,minimum width=1.65cm] (increment) at (0,5.65) {$k=k+1$};
  \node[decision] (test) at (0,4.05) {$a<\dfrac14$};
  \node[io,minimum width=1.90cm] (output) at (0,2.48) {输出 $k$};
  \node[terminal] (finish) at (0,1.18) {结束};

  % Main flow and the single vertical entry into a=aq.
  \coordinate (loopjoin) at (0,7.55);
  \coordinate (loopleft) at (-2.65,4.05);
  \draw[->] (start.south) -- (init.north);
  \draw[->] (init.south) -- (loopjoin) -- (multiply.north);
  \draw[->] (multiply.south) -- (increment.north);
  \draw[->] (increment.south) -- (test.north);
  \draw[->] (test.south) -- node[flowtext,right] {是} (output.north);
  \draw[->] (output.south) -- (finish.north);

  % 否 returns left, then the horizontal arrow joins the main stem above a=aq.
  \draw[->] (test.west) -- node[flowtext,above] {否} (loopleft) |- (loopjoin);
\end{tikzpicture}
